\documentclass[letterpaper,11pt]{article}
\usepackage[T1]{fontenc}
\usepackage[utf8]{inputenc}
\usepackage{lmodern}
\usepackage[french]{babel}
\usepackage[margin=1in,headheight=28pt]{geometry}
\renewcommand{\familydefault}{\sfdefault}
\usepackage{amsmath,amssymb,amsthm,bm,enumitem}
\usepackage[dvipsnames]{xcolor}
\usepackage[small,compact]{titlesec}
\usepackage{fancyhdr,fancyvrb}
\usepackage[colorlinks=true,allcolors=blue]{hyperref}
\pagestyle{fancy}
\fancyhead[L]{STT-2100 \\ Régression}
\fancyhead[R]{A.-A. Vallée \\ Université Laval}
\fancyfoot[C]{\thepage}
\setlength{\parindent}{0pt}
\setlength{\parskip}{5pt}
\fvset{fontsize=\small,frame=single,framesep=2mm}
\newcommand{\esp}{\mathbb{E}}
\newcommand{\var}{\operatorname{Var}}
\newcommand{\bo}[1]{\bm{#1}}
\newcommand{\N}{\mathcal{N}}
% Définir \avecsolutions avant l'input pour une version corrigée.
\newcommand{\prof}[1]{\ifdefined\avecsolutions\par{\color{blue}\textbf{Solution :} #1}\par\fi}
\begin{document}
\begin{center}{\Large Variables catégoriques et interactions avec R}\end{center}
Cette fiche accompagne les notes du chapitre 7. Les noms des objets peuvent être adaptés à vos données. Les exemples utilisent des données incluses dans R ou des données simulées; aucun fichier externe n’est requis.

\par\medskip\noindent\begin{minipage}{\linewidth}
\subsection*{Choisir le niveau de référence}
La fonction factor() traite une variable comme catégorique; relevel() choisit la référence. Avec les contrastes de traitement habituels, une variable à C catégories produit C-1 colonnes indicatrices en présence d'une constante.
\begin{Verbatim}
donnees <- iris
donnees$Species <- relevel(factor(donnees$Species),
                            ref = "setosa")
mod_add <- lm(Sepal.Length ~ Petal.Length + Species,
              data = donnees)
contrasts(donnees$Species)
head(model.matrix(mod_add))
summary(mod_add)
\end{Verbatim}

\end{minipage}\par
\par\medskip\noindent\begin{minipage}{\linewidth}
\subsection*{Inclure et tester les interactions}
L'opérateur * inclut les deux effets principaux et leur interaction. L'opérateur : inclut uniquement l'interaction. Un coefficient de Species représente ici une différence à Petal.Length=0, profil souvent peu réaliste.
\begin{Verbatim}
mod_int <- lm(Sepal.Length ~ Petal.Length * Species,
              data = donnees)
summary(mod_int)
# Test conjoint des deux interactions : pentes paralleles
anova(mod_add, mod_int, test = "F")
# Utilite globale de Species dans le modele additif
mod_sans <- lm(Sepal.Length ~ Petal.Length, data = donnees)
anova(mod_sans, mod_add, test = "F")
\end{Verbatim}

\end{minipage}\par
\par\medskip\noindent\begin{minipage}{\linewidth}
\subsection*{Centrer et interpréter les pentes}
Le centrage change la référence des coefficients sans changer les ajustements. Choisissez une référence pertinente et examinez le domaine de chaque espèce avant d'interpréter une différence entre groupes.
\begin{Verbatim}
donnees$petale_c <- donnees$Petal.Length - 4
mod_c <- lm(Sepal.Length ~ petale_c * Species,
            data = donnees)
summary(mod_c)
b <- coef(mod_int)
c(setosa = unname(b["Petal.Length"]),
  versicolor = unname(b["Petal.Length"] +
                       b["Petal.Length:Speciesversicolor"]),
  virginica = unname(b["Petal.Length"] +
                      b["Petal.Length:Speciesvirginica"]))
\end{Verbatim}

\end{minipage}\par
\par\medskip\noindent\begin{minipage}{\linewidth}
\subsection*{Prédire avec les mêmes niveaux}
Chaque profil ci-dessous utilise une longueur de pétale plausible pour son espèce. La variable catégorique de newdata doit reprendre les niveaux du modèle; une catégorie inconnue ne peut être prédite.
\begin{Verbatim}
profils <- data.frame(
  Petal.Length = c(1.5, 4, 5.5),
  Species = factor(c("setosa", "versicolor", "virginica"),
                   levels = levels(donnees$Species)))
predict(mod_int, newdata = profils, interval = "confidence")
\end{Verbatim}

\end{minipage}\par
\par\medskip\noindent\begin{minipage}{\linewidth}
\subsection*{Aide}
Dans R, utilisez \texttt{?lm} et l’aide des fonctions présentées. 
\end{minipage}\par
\end{document}
